\documentclass[10pt,a4paper]{amsart}
\usepackage[margin=19mm]{geometry}
\usepackage{amsmath,amssymb,mathtools}
\usepackage[scr=rsfs,cal=euler]{mathalfa}
\usepackage{xfrac}
\usepackage[unicode,hidelinks,bookmarks=false]{hyperref}
\newcommand{\ie}{i.e.\ }
\newcommand{\cf}{cf.\ }
\newcommand{\resp}{respectively\ }
\newcommand{\Subsection}{\subsection}
\providecommand{\nobreakdash}{\nobreak\hbox{-}}
\setlength{\emergencystretch}{2em}
\multlinegap0pt
\usepackage{amssymb}
\usepackage{dsfont}
\usepackage{mathtools}
\usepackage{xcolor}
\usepackage[shortlabels]{enumitem}
\newtheorem{theorem}{Theorem}
\newtheorem{proposition}{Proposition}
\title{On the Subtractive Ideal Structure of Commutative Semirings}
\author{Pubali Sengupta, Amartya Goswami, Pronay Biswas,, Sujit Kumar Sardar}
\date{Source: arXiv:2601.02120v2 (CC BY 4.0)}
\begin{document}
\maketitle

\noindent\emph{Source and licence.} On the Subtractive Ideal Structure of Commutative Semirings. Version arXiv:2601.02120v2 (CC BY 4.0), distributed by arXiv under the Creative Commons Attribution 4.0 International licence, http://creativecommons.org/licenses/by/4.0/, as recorded in the arXiv licence metadata for this exact version and shown on the abstract page https://arxiv.org/abs/2601.02120v2. Full licence notice: \texttt{licenses/arxiv-2601.02120-CC-BY-4.0.txt}.\par\smallskip

\noindent\textit{Editorial scope.} Counted unit: the article's proposition loc3 (source0.tex lines 735--745). The article's complete original proof of it is delivered with it (source0.tex lines 747--766, one \texttt{proof} environment of 20 lines). No mathematics is written, repaired or re-derived: the statement and the proof are the article's own lines, in the article's own notation, and no retained line is substituted or re-flowed.  Retained uncounted as the source-local context the proof needs: lines 382--384; lines 397--404.  Nothing was omitted from any retained span, so the package carries no omission record.  No retained line contains a citation, so the wrapper carries no bibliography and no bibliography entry.  No retained line contains a figure environment, an image-inclusion command or drawing code, so no image is delivered and the package declares no asset dependency.  Editorial formatting only, in addition: an AMS \texttt{amsart} preamble that restates the article's own declarations and macros used by the retained lines, namely the environments \texttt{theorem}, \texttt{proposition} on separate counters, together with the standard packages those lines need.  The retained environments print in this document as Theorem 1, Proposition 1 in order of appearance, whereas the article prints the same items with the numbering of its own full document.  The complete native source of the article as distributed in the arXiv e-print for this exact version is bundled unchanged as \texttt{sources/192134/source0.tex}.
\begin{theorem}\label{loc14}
    In a von Neumann regular semiring, the classes of $k$-prime and $k$-primary ideals coincide.
\end{theorem}

\begin{theorem}\label{loc11}
    Let $I$ be a $k$-ideal of a semiring $S$. Then the following conditions are equivalent.
    \begin{enumerate}
        \item $I$ is $k$-strongly irreducible.
        \item If $a$, $b \in S$ satisfying $\mathcal{C}_k(\langle a\rangle) \cap \mathcal{C}_k(\langle b\rangle) \subseteq I$, then $a \in I$ or $b \in I$.
        \item $S\backslash I$ is an $i_k$-system.
    \end{enumerate}
\end{theorem}

\begin{proposition}\label{loc3}
%Azizi, Theorem 2.1
Suppose $S$ is a semiring. 
\begin{enumerate}
\item If $I$ is a $k$-strongly irreducible ideal of $S$, then $I$ is a $k$-prime ideal of $S$ if and only if $I=\mathcal{R}_k(I)$.
\item If $S$ is a Laskerian semiring, then every $k$-strongly irreducible ideal is a $k$-primary ideal.

\item If $S$ is a von Neumann regular semiring, then an ideal is $k$-strongly irreducible if and only if
it is a $k$-primary ideal.
\end{enumerate}
\end{proposition}

\begin{proof}
	(1) If $I$ is $k$-prime, then clearly $I=\mathcal{R}_k(I)$. Conversely, suppose $I=\mathcal{R}_k(I)$. Let $AB \subseteq I$, where $A$, $B$ are two $k$-ideals of $S$. Then \[A\cap B \subseteq \mathcal{R}_k(A\cap B)=\mathcal{R}_k(A)\cap \mathcal{R}_k(B)=\mathcal{R}_k(AB)\subseteq \mathcal{R}_k(I)=I.\] Therefore, $A \subseteq I$ or $B \subseteq I$, as $I$ is $k$-strongly irreducible. Hence $I$ is $k$-prime.
    
        (2) Let $I$ be a $k$-strongly irreducible ideal of $S$ and suppose $\bigcap\limits_{i=1}^n A_i$ be a $k$-primary decomposition  of $I$. Then $\bigcap\limits_{i=1}^n A_i \subseteq I$. As $I$ is $k$-strongly irreducible, we have $A
        _j \subseteq I$, for some $j\in \{1,2,...,n\}$. This implies $A_j \subseteq I \subseteq A_j$, which further implies $I=A_j$. Thus $I$ is $k$-primary.
        
        (3) In a von Neumann regular semiring, $k$-prime and $k$-primary ideals coincide by Theorem~\ref{loc14}. Let $I$ be $k$-strongly irreducible and suppose $ab\in I$. It is enough, by Theorem~\ref{loc11}, to prove
\[
\mathcal C_k(\langle a\rangle)\cap\mathcal C_k(\langle b\rangle)\subseteq I.
\]
Take $c$ in the intersection. Choose $\alpha_1,\alpha_2,\beta_1,\beta_2\in S$ such that
\[
c+a\alpha_1=a\alpha_2,\qquad c+b\beta_1=b\beta_2.
\]
Put $K=\mathcal C_k(\langle ab\rangle)$. Multiplying the first equality by $b\beta_1$ and the second by $a\alpha_1$ shows, respectively, that $cb\beta_1\in K$ and $ca\alpha_1\in K$. Now multiplying the above equalities gives
\[
c^2+cb\beta_1+ca\alpha_1+ab\alpha_1\beta_1=ab\alpha_2\beta_2.
\]
By subtractivity therefore we get $c^2\in K$. Since $S$ is von Neumann regular, $c=c^2r$ for some $r\in S$, so $c\in K\subseteq I$. Thus Theorem~\ref{loc11} gives $a\in I$ or $b\in I$, and $I$ is $k$-prime, hence $k$-primary. Conversely, every $k$-primary ideal is $k$-prime by Theorem~\ref{loc14}, and every $k$-prime ideal is $k$-strongly irreducible.
\end{proof}

\end{document}
